\hypertarget{appendix-8f6a01f284a7}{}\label{referencestart:appendix-8f6a01f284a7}
\section{Conditioning and singularity}\label{appendix-8f6a01f284a7-conditioning-and-singularity}

\hypertarget{appendix-8f6a01f284a7-conditioning-and-singularity}{}

\textbf{Curriculum depth:} Developed · \textbf{First seen:} C004 · \textbf{Latest development:} Phase7

\begin{allmightroledefinition}

\subsection{Definition and mechanism}\label{appendix-8f6a01f284a7-definition-and-mechanism}

\hypertarget{appendix-8f6a01f284a7-definition-and-mechanism}{}

Conditioning measures sensitivity to input perturbation. \(A=10^{-4}I_2\) has determinant \(10^{-8}\) but condition number one in the 2-norm. Uniform scaling changes determinant magnitude without changing this conditioning measure.

\end{allmightroledefinition}

Phase 7 develops \(\kappa(A)=\|A\|\|A^{-1}\|\), near-parallel-column sensitivity, pivoted elimination and inverse residuals. Equal determinants can conceal very different stretch ratios: \(\operatorname{diag}(10^{-10},10^{-10})\) and \(\operatorname{diag}(1,10^{-20})\) both have determinant \(10^{-20}\) but conditions one and \(10^{20}\). \href{../02-Mathematics-for-Visual-Computing-P01/main.pdf\#nameddest=01-foundations-i-ch018-an-inverse-can-exist-and-still-amplify-uncertainty}{Full worked sensitivity example (Mathematics for Visual Computing, Part 1, Ch. 2)}.

Subtraction of nearby inputs exposes input uncertainty relative to a small difference. Sensitivity belongs to the problem; exact subtraction of stored operands cannot recover information lost when those inputs were rounded.

\subsection{Prerequisites}\label{appendix-8f6a01f284a7-prerequisites}

\hypertarget{appendix-8f6a01f284a7-prerequisites}{}

\begin{itemize}
\tightlist
\item
  \hyperref[appendix-0ec599f8ba27]{Numerical tolerance policies} --- An absolute numerical threshold can change with scale without measuring mathematical sensitivity.
\end{itemize}

\subsection{Used by}\label{appendix-8f6a01f284a7-used-by}

\hypertarget{appendix-8f6a01f284a7-used-by}{}

\begin{itemize}
\item
  \href{../02-Mathematics-for-Visual-Computing-P01/main.pdf\#nameddest=appendix-48d2e66378ab}{Matrix inversion pivoting and residuals (Mathematics for Visual Computing, Part 1, Reference notes: Matrix inversion pivoting and residuals)}
\item
  \hyperref[appendix-a7fbaf6d1b1c]{Forward and backward error} --- Input sensitivity is distinct from errors an implementation introduces.
\item
  \hyperref[appendix-a2773a6a2f9d]{Cancellation and loss of significance} --- Input sensitivity is distinct from errors an implementation introduces.
\end{itemize}

Connect later dependent concepts here as their lessons arrive. - \href{../07-Vision-XR-and-Playable-Worlds-P01/main.pdf\#nameddest=appendix-f6e77f4a200f}{Gaussian elimination and pivoting (Vision, XR, and Playable Worlds, Part 1, Reference notes: Gaussian elimination and pivoting)} --- Extends this foundation into linear-systems and decomposition analysis. - \href{../07-Vision-XR-and-Playable-Worlds-P01/main.pdf\#nameddest=appendix-c3421e6b4f2c}{QR factorization and least squares (Vision, XR, and Playable Worlds, Part 1, Reference notes: QR factorization and least squares)} --- Extends this foundation into linear-systems and decomposition analysis. - \href{../04-Light-Transport-and-Image-Synthesis-P01/main.pdf\#nameddest=appendix-18eacf61a04b}{ODE integration and stability (Light Transport and Image Synthesis, Part 1, Reference notes: ODE integration and stability)}

\begin{itemize}
\tightlist
\item
  \href{../07-Vision-XR-and-Playable-Worlds-P01/main.pdf\#nameddest=appendix-d83b0ad1be7f}{A posteriori error bounds for contractive iteration (Vision, XR, and Playable Worlds, Part 1, Reference notes: A posteriori error bounds for contractive iteration)}
\item
  \href{../07-Vision-XR-and-Playable-Worlds-P01/main.pdf\#nameddest=appendix-aa2bf0d79639}{Algorithmic stability (Vision, XR, and Playable Worlds, Part 1, Reference notes: Algorithmic stability)}
\item
  \href{../07-Vision-XR-and-Playable-Worlds-P01/main.pdf\#nameddest=appendix-f63eb850796d}{Bracketed and open root solvers (Vision, XR, and Playable Worlds, Part 1, Reference notes: Bracketed and open root solvers)}
\item
  \href{../07-Vision-XR-and-Playable-Worlds-P01/main.pdf\#nameddest=appendix-31b8c295e05f}{Numerical error budgets (Vision, XR, and Playable Worlds, Part 1, Reference notes: Numerical error budgets)}
\item
  \href{../07-Vision-XR-and-Playable-Worlds-P01/main.pdf\#nameddest=appendix-b9afba7f63e5}{Stable quadratic root evaluation (Vision, XR, and Playable Worlds, Part 1, Reference notes: Stable quadratic root evaluation)} \#\# Recall and next depth
\end{itemize}

Reconstruct the concrete example in the definition and identify its validity assumptions.

\textbf{Pending:} Robust condition estimation, growth factors, advanced factorizations and application-specific error budgets extend the current inversion laboratory.

\subsection{Source and navigation}\label{appendix-8f6a01f284a7-source-and-navigation}

\hypertarget{appendix-8f6a01f284a7-source-and-navigation}{}

\hyperref[01-foundations-i-ch001-when-arithmetic-meets-uncertainty]{M005 development}

\hyperref[01-foundations-i-ch001-a-stretching-ruler-for-real-valued-quantities]{M004 teaching chapter} · \hyperref[appendix-fd090d982a15-m004---parts-61-76]{Per-part definitions} · \href{../07-Vision-XR-and-Playable-Worlds-P01/main.pdf\#nameddest=appendix-4ad47bdd269a}{Dependency Map (Vision, XR, and Playable Worlds, Part 1, Reference notes: Dependency Map)}
